\chapter{Discussion}

Previous models of adult neurogenesis and pattern separation have been confined to the DG~\parencite{chavlisDendrites2017,
kimEffect2024, yangDynamic2025}, leaving open whether the circuit-wide sparsification observed \textit{in
vivo}~\parencite{mchughAdultborn2022} follows from the known DG-CA3 microcircuitry. In the model built here, the DG reproduces the
division of roles reported by~\textcite{kimEffect2024}, the mGCs separating their inputs while the iGCs integrate them, with
separation improving through sparser coding~\parencite{chavlisDendrites2017}. The same computation, together with the interneuron
recruitment it drives, extends to CA3, whose activity also becomes sparser, consistent with a modulatory rather than an encoding
role for iGCs~\parencite{mchughAdultborn2022, folszAdultborn2023}.

The gain in mGC separation follows from a redistribution of activity within the GC population. Increasing the EC$\to$iGC
connectivity raises the iGC activation, and the iGCs drive the BCs and the HIPPs both directly and through the MCs; the resulting
inhibition lowers the mGC activation, and the sparser mGC output raises its separation degree. Since the total GC activation stays
close to the control value at every connectivity level, the DG does not become less active, and the activity given up by the mGCs
is taken up by the iGCs. The separation degree measured over the whole GC population moves in the opposite direction for the same
reason: the iGCs are few and are active for a large share of the patterns, so they are shared between patterns, and pooling them
with the mGCs adds correlation to the output. That decline is accounted for almost entirely by the loss of decorrelation, the
activation term showing no trend across connectivity levels.\ \textcite{kimEffect2024} varied the same quantity and found two
states divided by a threshold, the separation degree falling below that of a network without iGCs at low synaptic maturity, where
the reduced innervation of the iGCs lowers their activation enough to withdraw inhibition from the mGCs. No such state appears
here, the mGC separation degree staying above the control at every level. The difference is consistent with how the two models
reduce the input to the iGCs:~\textcite{kimEffect2024} scale the connection probability from the EC and the MCs together, whereas
here only the EC$\to$iGC probability is scaled and the MC$\to$iGC projection is left intact, so the iGCs retain enough excitatory
drive to engage the BCs and the HIPPs even at the lowest EC$\to$iGC connectivity.

The pCA3 separation degree rises with neurogenesis but stays below one. Of the two factors combined in Equation~\ref{eq:dp}, it is
the activation term that holds it there: the pCA3 population is active at 24\% to 30\% against the 10\% of the input patterns, and
the measure returns values below one for any code denser than its input, whichever way the correlations move. Measured against the
GC output that drives it rather than against the EC input, the pCA3 code is the more correlated of the two at low EC$\to$iGC
excitability and the less correlated at high ones, its own output correlation staying between $0.22$ and $0.28$ while that of the
GC patterns rises from $0.14$ to $0.55$. The balance between the two factors also differs from the one found in the DG, the rise
in the pCA3 separation degree being divided roughly two to one between the fall in activation and the gain in orthogonalization,
against a much larger share for activation in the mGCs.

In agreement with~\textcite{mchughAdultborn2022}, activating the iGCs recruited inhibitory interneurons in both the DG and CA3, the
circuit counterpart of the increased GC population sparsity of Section~\ref{sec:res_sparsity}. This recruitment was not a uniform
elevation of firing: of all the populations, only the BCs show a spectral peak standing clearly above the background of their own
spectrum (Figure~\ref{fig:frequency_spectra}), so the perisomatic inhibition engaged by the pulse is organized as a population
rhythm.
The resolution set by the analysis window locates that frequency only to within one bin, so it should be taken as approximate. Two
features of the model response differ from the \textit{in vivo} findings: the mGCs showed a net excitation rather than the
suppression expected of a principal cell, and the pCA3 cells showed a small inhibition, for which~\textcite{mchughAdultborn2022}
reported no effect. The area-level recordings of that study could not separate the inhibitory contributions of the HIPPs and the
BCs, so the divergence between the two DG interneuron populations cannot be adjudicated against their data.

Neither the direct iGC$\to$mGC inhibitory synapse nor the feedforward and feedback inhibition triggered by the pulse suppressed
the mGCs. This residual excitation suggests that the iGC$\to$mGC interaction is captured only in part, and that an inhibitory
contribution is missing, plausibly from interneuron types absent from the reduced circuit: at least 13 morpho-physiological
interneuron types have been distinguished in the DG~\parencite{degroInterneuron2022}, against the two represented here.

The size of the effect produced by the pulse depends on how much the iGCs already contribute at baseline. The same current is
delivered to the same number of iGCs at every level and adds a roughly constant increment to the iGC population rate, while the
baseline iGC rate rises about elevenfold across the connectivity range, so the pulse raises that rate some thirtyfold at the
lowest connectivity and about fourfold at the highest. The gain in GC population sparsity produced by the pulse falls over the
same range (Figure~\ref{fig:opto_sparsity}). The sparsification that follows a brief activation of adult-born GCs would therefore be
expected to depend on how far the activated cells have integrated into the circuit, and to be largest for the least mature of
them.

Several simplifications bound these results. Each neuron is reduced to a single compartment, so the dendritic processing that
controls GC sparsity and separation~\parencite{chavlisDendrites2017}, and that varies across GC subtypes~\parencite{maoExcitatory2025},
is not represented. The interneuron repertoire is minimal, comprising a single BC and HIPP type in the DG and one inhibitory type
in CA3~\parencite{kopsickFormation2024}, a likely source of the unresolved mGC response. The iGCs are represented as a single static
population, so the sweep over EC$\to$iGC connectivity samples the maturation range rather than following cells through it.

The effects on CA3 require further investigation. Because its synapses are held fixed here, the model cannot address how adult
neurogenesis affects storage and recall in the recurrent CA3 network; introducing synaptic plasticity would allow this to be
studied, along with the influence of adult-born cells on the formation and retrieval of CA3 engrams~\parencite{rollsMechanisms2013,
kopsickFormation2024, koSystems2025}.
